\documentclass[dvipdfmx]{jarticle}
\pagestyle{empty}
\newcommand{\ttsc}{\texttt{;}}
\newcommand{\ttcol}{\texttt{:}}
\newcommand{\tteq}{\texttt{=}}
\newcommand{\ttasgn}{\ttcol\tteq}
\newcommand{\ttpl}{\texttt{+}}
\newcommand{\ttbrkl}{\texttt{[}}
\newcommand{\ttbrkr}{\texttt{]}}
\begin{document}
\begin{tabular}{|l|l|}\hline
生成規則 & 型制約規則\\\hline
$E\rightarrow c$ & $E.\mathsf{type}\ttasgn\mathsf{char}$\\
$E\rightarrow n$ & $E.\mathsf{type}\ttasgn\mathsf{integer}$\\
$E\rightarrow i$ & $E.\mathsf{type}\ttasgn i.\mathsf{type}$ (記号表を検索する)\\
$E\rightarrow E_1\ttpl E_2$ & $E.\mathsf{type}\ttasgn\ \mathbf{if}\ 
E_1.\mathsf{type}=\mathsf{integer}\ \mathbf{and}\ E_2.\mathsf{type}=\mathsf{integer}
\ \mathbf{then}\ \mathsf{integer}\ \mathbf{else}\ \mathsf{type\_error}$\\
$E\rightarrow \texttt{*}E_1$ & $E.\mathsf{type}\ttasgn\ \mathbf{if}\ E_1.\mathsf{type}=\mathsf{pointer}(t)
\ \mathbf{then}\ t\ \mathbf{else}\ \mathsf{type\_error}$\\
$E\rightarrow E_1\ttbrkl E_2\ttbrkr$ & $E.\mathsf{type}\ttasgn\
\mathbf{if}\ E_1.\mathsf{type}=\mathsf{array}(s,t)\ \mathbf{and}\ E_2.\mathsf{type}=\mathsf{integer}
\ \mathbf{else}\ \mathsf{type\_error}$\\
$S\rightarrow i\mathtt{:=} E$ & $i.\mathsf{type}\!=E.\mathsf{type}$\\\hline
\end{tabular}
\end{document}